\section{All arcs through a double contact}
\label{sec:all-arcs-v166}
The six coefficient directions of $B_2$ need not vanish in pairs,
and a collision need not have a chosen square root. In this section
we give a finite presentation of the full pullback, its horizontal
closure, and its vertical torsion for every discrete valuation ring
arc through $o=(x^2,0,0)$. The presentation is by a fixed collection
of coordinate systems with their recovery opens. It is not a list
of irreducible components indexed by the valuation of a single
chosen coefficient.

Let $R$ be a discrete valuation ring which is a $\C$-algebra, with
uniformizer $\tau$ and residue field $\C$. The same statements hold
after extending that residue field. Write an arbitrary arc centred
at $o$ uniquely as
\begin{equation}\label{eq:all-arc-coefficients-v166}
 f=(x-c)^2+d,\qquad g=u(x-c)+v,\qquad
 r=w(x-c)+z_0,
\end{equation}
where $c,d,u,v,w,z_0\in\tau R$. No relation between these six
elements is assumed. In particular neither $g=r=0$ nor a prescribed
ramification order is imposed. Set
\[
 H_R=\widetilde\Gamma_2\times_{B_2}\operatorname{Spec}R,
 \qquad Y_R=\overline{(H_R)_K}^{\,H_R},
\]
where $K=\operatorname{Frac}R$ and total-graph normalization precedes
the base change.

\subsection{Five conic coordinates and two division coordinates}
The vector space of conic equations through $P=[1:0:0]$ has basis
$FG,FR,G^2,GR,R^2$. Evaluation at the following five vectors gives
linearly independent functionals:
\begin{equation}\label{eq:five-vectors-v166}
 (0,0,1),\quad(0,1,0),\quad(0,1,1),\quad(1,1,0),\quad(1,0,1).
\end{equation}
Here $R$ in a target coordinate denotes the third homogeneous
coordinate, not the discrete valuation ring. To avoid confusion
below it will only occur in homogeneous equations.
Choose matrices $M_i$ whose first column is $(1,0,0)$ and whose
third columns are the vectors in \eqref{eq:five-vectors-v166}:
\begin{equation}\label{eq:five-matrices-v166}
\begin{gathered}
 M_0=I_3,\qquad
 M_1=\begin{pmatrix}1&0&0\\0&0&1\\0&1&0\end{pmatrix},\qquad
 M_2=\begin{pmatrix}1&0&0\\0&1&1\\0&0&1\end{pmatrix},\\
 M_3=\begin{pmatrix}1&0&1\\0&0&1\\0&1&0\end{pmatrix},\qquad
 M_4=\begin{pmatrix}1&0&1\\0&1&0\\0&0&1\end{pmatrix}.
\end{gathered}
\end{equation}
For coordinate system $i$ replace the coefficient column by
$M_i^{-1}(f,g,r)^{\mathsf t}$, and recompute the six coefficients
in \eqref{eq:all-arc-coefficients-v166}. The first polynomial stays
monic of degree two and the other two have degree at most one.
The central point remains $o$. These are changes of target
coordinates, always undone when the embedded curve is returned to
the fixed target; they are not quotient identifications of Hilbert
points.

For any one of these transformed six-tuples put
\begin{equation}\label{eq:all-arc-HB-v166}
\begin{split}
 D_i&=R[\lambda,e,A,B,C],\\
 K_i&=(e^2C-d,\ eB-u,\ e^2A+v,\
                 \lambda u-eA-w,\ \lambda v-e^2BC-z_0)\subset D_i.
\end{split}
\end{equation}
There are also two division coordinate systems. The first has
\begin{equation}\label{eq:all-arc-division-v166}
 D_+=R[\alpha,\beta],\qquad
 K_+=(\alpha u+\beta v-w,\ \alpha v-\beta ud-z_0).
\end{equation}
The second is obtained by interchanging the pairs $(u,v)$ and
$(w,z_0)$ and is denoted $(D_-,K_-)$.

\begin{theorem}[Finite-chart miniversal specialization]
\label{thm:all-arcs-v166}
For every arc \eqref{eq:all-arc-coefficients-v166}, the seven
coordinate systems \eqref{eq:all-arc-HB-v166} and
\eqref{eq:all-arc-division-v166}, restricted to their Hilbert recovery
opens, cover $H_R$. On each such open its ring is the corresponding
$D_\bullet/K_\bullet$, its horizontal ring is
\begin{equation}\label{eq:all-arc-horizontal-v166}
 D_\bullet/(K_\bullet:\tau^\infty),
\end{equation}
and its vertical torsion ideal is
$(K_\bullet:\tau^\infty)/K_\bullet$. These formulas retain the
scheme structures, including all embedded primes and nilpotents.
They commute with localization, completion, and flat extensions of
discrete valuation rings. They determine the entire specialization,
not merely a neighbourhood of one chosen conic or primitive tail.

The recovery opens are the opens of the total coefficient graph,
not extra conditions on the six coefficients of the arc. Thus the
theorem permits nonzero $g$ and $r$ jets, arbitrary relative orders
of vanishing, unequal residual directions, and arbitrary ramification.
It includes arcs entirely contained in an incidence stratum.
\end{theorem}
\begin{proof}
In the Hilbert--Burch coordinates of
Lemma~\ref{lem:HB-family-v163}, put $z=x-c$ and $H=R-\lambda G$.
The base map reads
\[
 f=z^2+e^2C,\quad g=eBz-e^2A,\quad
 r=\lambda g-eAz-e^2BC.
\]
Equating its coefficients with those of the retained arc gives
exactly the five equations in \eqref{eq:all-arc-HB-v166}; the sixth
coefficient eliminates $c$. In particular the constant equation
$\lambda v-e^2BC=z_0$ must not be discarded for a general arc.
No radical has been taken. The total chart is regular on its
recovery open, so the same computation applies to the normalization
of the total graph.

In the primitive chart write the recovered direction as
$\alpha+\beta z$. Division by $f=z^2+d$ gives
\[
 (\alpha+\beta z)(uz+v)-\beta u(z^2+d)
       =(\alpha u+\beta v)z+\alpha v-\beta ud.
\]
The condition that this equal $r$ proves
\eqref{eq:all-arc-division-v166}. Interchanging $g$ and $r$ gives
the other primitive chart. These are the degree-two division charts
of Theorem~\ref{thm:division-chart-v162}, after the invertible change
between $x$ and $z$ coordinates. Their equations also retain the
coefficient projection.

We specify the opens used here. On a conic chart the restriction maps
in degrees $(0,2)$ and $(1,1)$ have ranks five and four. Restrict to
the opens defined by these ranks, by the nonzero selected conic
coefficient, and by the $tG,tH$ pivot for the two bilinear equations.
The normalized inverse is precisely the recovery map in
Lemma~\ref{lem:HB-family-v163}; the matrices in
\eqref{eq:five-matrices-v166} transport these opens. They contain
every central conic on which the corresponding evaluation functional
is nonzero. On a division chart use the split-injection open in
Lemma~\ref{lem:recovery-details-v163}. These opens contain every
central primitive tail for which the selected homogeneous direction
coordinate is a unit. Localizing the displayed rings to these opens
is understood; no global affine-isomorphism claim outside the
recovery domains is needed.

A nonzero conic through $P$ cannot vanish on all five vectors in
\eqref{eq:five-vectors-v166}. Indeed their evaluations recover,
in order, the coefficients of $R^2$, $G^2$, $GR$, $FG$, and $FR$.
Thus the five conic coordinate systems cover the complete conic
component. A primitive jet over $\C[x]/(x^2)$ is represented by a
pair with at least one unit, and therefore belongs to one of the two
division coordinate systems. The complete-fibre theorem
\ref{thm:complete-fibre-v163} shows that these are all central
points, including the doubled-line intersection. The same charts
compute their scheme structures, not only their closed points.

Their union is an open neighbourhood of the entire special fibre
in $H_R$. Its complement is closed in the proper $R$-scheme $H_R$.
A nonempty closed subset of a proper scheme over a local discrete
valuation ring has image meeting its closed point. Since the
complement misses that fibre, it is empty. This proves exhaustion
of $H_R$, including its generic fibre. Now apply
Proposition~\ref{prop:dvr-horizontal-v166} on the seven coordinate
domains. Its uniqueness assertion glues the horizontal quotients.
Lemma~\ref{lem:saturation-v166} proves compatibility on overlaps
and with the stated flat changes of rings. The torsion formula
follows from the same proposition.
\end{proof}

\begin{corollary}[Universal interpretation of the seven presentations]
\label{cor:all-arcs-universal-v166}
Choose the stratum containing the generic point of the arc in
Theorem~\ref{thm:all-arcs-v166}. Its seven horizontal presentations
are the pullback of the universal family on the corresponding
modification in Theorem~\ref{thm:stratum-flattening-v166}. Hence
specialization does not require an additional choice of a component,
a normalization of the special fibre, or a stable-map replacement.
\end{corollary}
\begin{proof}
Both constructions are flat embedded quotients with the original
generic fibre. Uniqueness in
Proposition~\ref{prop:dvr-horizontal-v166} identifies them on every
coordinate domain and hence globally.
\end{proof}

The adjective miniversal here refers to the full six-dimensional
coefficient base and its framed fixed-target comparison. It does
not assert a finite list of valuation chambers, nor a classification
of the full fibre for $a\geq3$. The theorem gives the exact ideals
and their saturation for every arc through the double contact. Its
specialization to $g=r=0$ recovers the previously computed collision;
it is not restricted to that subbase.

\subsection{Which conics are horizontal?}
The vertical component in the symmetric collision has a concrete
modular explanation: it records limits approached in coefficient
directions other than the double-incidence subbase. Let
\[
 \mathcal Q_P=\Pj\langle FG,FR,G^2,GR,R^2\rangle,\qquad
 S=V(q_{FG},q_{FR})\subset\mathcal Q_P.
\]
Thus $S\simeq\Pj^2$ is the linear space of conics singular at $P$.
Its ideal is generated by the two nontrivial first derivatives at
$P$; this is an intrinsic first-jet condition on the conic equation.

\begin{theorem}[Arc realization and the meaning of vertical excess]
\label{thm:conic-realization-v166}
Every conic $Q\in\mathcal Q_P$, including a doubled line, occurs as
the tail of the special Hilbert curve of an arc in $B_2$ whose
generic point is nonincident. Among arcs in the double-incidence
subbase $g=r=0$ with $f$ generically squarefree, the conic tails
reached horizontally are exactly $S$.
\end{theorem}
\begin{proof}
Choose one of the five coordinates so that $Q$ is
\[
 (R-\lambda_0G)^2+B_0F(R-\lambda_0G)+A_0FG+C_0G^2=0.
\]
Take $e=\tau$, $c=0$, $\lambda=\lambda_0$, and choose polynomial
series $A(\tau),B(\tau),C(\tau)$ with prescribed constant terms
$A_0,B_0,C_0$ and
$A(\tau)^2+C(\tau)B(\tau)^2\ne0$ in $K$. Such a choice always
exists because this polynomial is not identically zero on the
three-dimensional coefficient space; a generic line through the
specified point is not contained in its zero locus. The
Hilbert--Burch base map produces an arc centred at $o$. On its
generic point $e(A^2+CB^2)\ne0$, so the three quadratics span the
space of binary quadratics and define the nonincident graph. The
flat determinantal family has special curve $C_0\cup Q$ by
Lemma~\ref{lem:HB-family-v163}. Undoing the constant coordinate
change gives the required embedded arc in the original target.

For the second assertion use any conic coordinate system for a
horizontal limit; the transformations in
\eqref{eq:five-matrices-v166} preserve the condition $g=r=0$.
On the generic fibre the coefficient equations give $eB=eA=0$.
Also $e\ne0$ there, since $e^2C=d$ and the centred quadratic is
generically squarefree, so $d\ne0$. Thus $A=B=0$ on the horizontal
closure. Its conic limit lies in $S$. Conversely, in a chart for
any $Q\in S$, set $A=B=0$, $e=\tau$, and take $C(\tau)$ with the
specified constant term and nonzero generic value. Then
$f=x^2+\tau^2C(\tau)$ is generically squarefree and $g=r=0$.
The displayed flat Hilbert family gives a point of the horizontal
closure: it agrees with that closure generically and factors through
its saturated ideal over the torsion-free ring $R$. Its conic limit
is $Q$. This includes doubled lines by taking $C(0)=0$ and
$C(\tau)\ne0$.
\end{proof}

The first assertion concerns closures of individual generic Hilbert
points on generically nonincident arcs. The second concerns points
of the horizontal family of complete incidence fibres. The two
notions agree on their common incidence interpretation but should
not be conflated. In particular the vertical $\Pj^4$ is not a
collection of spurious curves: its complement of $S$ consists of
valid limits that require different coefficient directions. No
variation of a stability parameter is involved.
